\begin{table}[H]
	\centering
	\fontsize{7}{7}\selectfont
	\setlength{\tabcolsep}{3pt}
	\renewcommand{\arraystretch}{1.4}
	\begin{tabularx}{\textwidth}{%
		>{\raggedright\arraybackslash}p{1.7cm}  
		>{\raggedright\arraybackslash}p{1.5cm}    
		>{\raggedright\arraybackslash}p{2.3cm}  
		>{\raggedright\arraybackslash}p{1.4cm}  
		>{\raggedright\arraybackslash}p{2.0cm}  
		>{\raggedright\arraybackslash}p{2.5cm}    
		>{\raggedright\arraybackslash}p{2.0cm}} 
		\toprule
		\textbf{Interesado} & \textbf{Rol} & \textbf{Relación} & \textbf{Impacto} & \textbf{Poder de influencia} & \textbf{Interés} & \textbf{Compromiso} \\
		\midrule
		Docentes de Ingeniería de Sistemas & Usuarios clave & Utilizarán los modelos para apoyar la docencia en arquitectura empresarial e Infraestructura de~\TI & Medio-Alto & Medio, pueden aportar sugerencias metodológicas o pedagógicas & Alto, buscan fortalecer la enseñanza con herramientas de modelado y representación visual & Medio-Alto, dependiendo de la integración en sus asignaturas \\
		\midrule
		Estudiantes de Ingeniería de Sistemas & Beneficiarios directos & Usarán los modelos y herramientas como apoyo en proyectos académicos y de investigación & Medio & Bajo, sin capacidad de decisión pero con alto valor en la adopción y validación práctica & Medio-Alto, el uso de notaciones de modelado mejora sus competencias técnicas y analíticas & Medio, si perciben utilidad práctica en su formación \\
		\midrule
		Grupo de Investigación GRID & Beneficiario principal & Lidera la iniciativa de modelar la Infraestructura de~\TI mediante lenguajes de arquitectura empresarial & Alto & Alto, autoridad en la definición y uso de herramientas dentro del proyecto & Alto, busca mejorar la gestión del conocimiento y planificación tecnológica del grupo & Alto, ya que los resultados impactan directamente sus procesos y capacidades \\
		\midrule
		Facultad de Ingeniería & Facilitador institucional & Brinda respaldo académico e institucional al proyecto, garantizando recursos y visibilidad & Alto & Alto, con poder para integrar los resultados en otros programas o proyectos institucionales & Medio, enfocado en la alineación con objetivos estratégicos y académicos & Medio, especialmente si el proyecto demuestra valor institucional \\
		\midrule
		Grupos de investigación aliados & Colaboradores técnicos & Podrían replicar o adaptar los modelos para sus propias infraestructuras o líneas de investigación & Medio & Medio-Bajo, influyen mediante colaboración y retroalimentación técnica & Alto, interesados en compartir metodologías y resultados reproducibles & Medio-Alto, según su participación activa en la cooperación \\
		\midrule
		Comunidad académica y profesional de arquitectura empresarial & Proveedores y consumidores de conocimiento & Aportan buenas prácticas y frameworks estandarizados, y se benefician de las experiencias documentadas por el grupo & Medio-Alto & Bajo, no inciden directamente en decisiones internas del proyecto & Medio, la adopción de ArchiMate y otras notaciones fortalece la comunidad académica & Bajo-Medio, su compromiso depende de la difusión de resultados y publicaciones \\
		\bottomrule
	\end{tabularx}
	\caption{Análisis de \textit{stakeholders} del proyecto de modelado de Infraestructura de~\TI}
	\label{tab:stakeholders}
\end{table}
